\section{Limitations}\label{sec:limits}
\textbf{Measurement.} The counter is uncalibrated; it was compared with \texttt{micros()} up to \EspHiLongest\,s on one weak point and up to \RatioLongSEspHi--\RatioLongSEspLo\,s on non-weak ESP points; its representable limit $T_{\max}$ = \TMaxS\,s (\TMaxDays\ days) is not a guarantee of correctness. A wrap artefact (\WrapS\,s) is removed by R5, which was set after the data were seen. Counter pairing failed on most long \espA\ Phase~2 points, on \espB\ \grover\ at $n=12$ and on several Phase~3 points of all boards (cause unverified; the Uno points at the memory wall never ran); their \texttt{micros()} data are used, only the agreement lacks them. The Uno timer resolves \UnoResUs\us, the ESP timers \EspResUs\us; the models ignore points below \ModelFloorUs\us. The loopback self-test offset of \SelfTestLo--\SelfTestHi\us\ cannot be split between generator and counter, and the target input pad of the counter was not tested. The unwrapping of the counter timestamps needs both edge stamps to lie within about 27\,s of the true width. Rule R5 changed \RFivePoints\ points. In the offload decomposition the residual absorbs the error of subtracting two medians, so it can be negative at large $n$ and is bookkeeping, not a measurement.

\textbf{Scope.} One unit per board, one PC, one baud rate (\BaudSp); the crossover was not tested elsewhere. Boards return timing only: there are no board-level correctness data, and the microcontroller \grover\ builds may or may not lose norm like the PC build. Phase~3 samples one shot per request and compares it with nothing. \ProvDirty\ of \NSessions\ sessions carry a dirty-tree flag; \ProvCommitted\ have a firmware source hash equal to a committed tree and \ProvUncommitted\ ran uncommitted firmware; no hash of the flashed binary was recorded.

\textbf{Weak and mixed evidence.} Phase~2 \grover\ points at $n=\WeakPtwoEspHi$ (\espB) and $n=\WeakPtwoEspLo$ (\espA) and Phase~3 \grover\ points at $n=\WeakPthreeEspHi$ (\espB) have fewer than \WeakMin\ independent repetitions or requests and are never primary evidence. All selected \espA\ \grover\ points are firmware class~B (the timing relevance of the class is read from the firmware diff: interpretation) and so outside the fit and the clean scores; \grover\ on the ESP boards confounds firmware version and repetition count, so it is not a model result; \espB\ \grover\ comes from two hashes, both class~A by reading the diff (unverified: no binary hash).

\textbf{Post-hoc choices (all set after data or scores were seen).} The weak rule (first \WeakFirst, then \WeakMin\ repetitions after the first scores), the firmware classes, the identifiability rule and its boundary condition, the effective-repetition check, the primary Phase~1 criterion (infidelity of the normalised state, chosen after the loss of norm was identified as the cause of the raw failure of \grover), the session selection and the counter-session fallback (chosen because those sessions have usable pairs, an availability-conditioned choice), rule R5, the choice of which $B$ predictor to use per board (made with the walls known), the band-and-bracket design of the crossover, and, as new views of the same data, the per-circuit crossover figure, the band check (\BandTolPct\,\% tolerance, one band per board), the decision map, the predicted-crossover comparison (with one residual constant per board and the \SimSmallMs\,ms cut that defines it), the fast-versus-slow shift counts, the ratio and breakdown views (with the \BreakPcShare\,\% threshold) and the choice of the decision-map classes. The band check could only fire: one band per board was the only drawing implemented. Also written after the data were collected: the counter rules R1--R4 and the forms of M1 and M2 (the latter before the first scores, per the authors' record). The held-out-size sensitivity was added after the first scores. Fixed before scoring, by the authors' record and not verifiable from the commit history: the split (largest three $n$) and the \ModelFloorUs\us\ floor. M2 was not defined blind with respect to \grover. The \espA\ M2 error moved from \PrevMtwoEspLo\,\% in an earlier analysis round of this study to \MapeTwoEspLo\,\% after these rules, and the \espB\ M1 error from \PrevMoneEspHi\,\% to \MapeOneEspHi\,\%, without a split by rule. The condition-number cap (\CondCap) never binds (largest value \CondMaxTwo); the rule that binds is the zero-boundary condition, which was not set in advance. The effect of the weak switch alone on M2 is \WeakSwitchWithout\,\% against \WeakSwitchWith\,\%; the effect of the other rules and of the session selection together was not split.

\textbf{Interpretation.} The fitted constants are not physical gate costs; the \espA\ CNOT cost is impossible as a per-iteration cost. The \SelfTestLo--\SelfTestHi\us\ offset, the Uno timebase explanation, the cause of the PC speed change and the repeated rows in \espB\ \grover\ are interpretation or unverified and not findings. The memory test resolves $B$ only to a factor of two and the Uno confirmation is close to circular.
